\documentclass[12pt]{article}
\usepackage[utf8]{inputenc}
\usepackage[T1]{fontenc}
\usepackage{hyperref}
\title{Southern Kalapuya Language Resources for K-12 Curriculum Development}
\author{Compiled by Claude Sonnet 5.5 (claude-sonnet-5-5) for Andr\'as Kornai}
\date{}
\begin{document}
\maketitle

\section{Introduction}
Southern Kalapuya (Yoncalla) is a Kalapuyan language of the Willamette and Umpqua valleys in Oregon. A Wikipedia entry on Yoncalla lists it as extinct, and gives Laura Albertson as the last known partial speaker, attested in 1937. Sources disagree on later dates. Henry Zenk is reported to argue that Albertson, who died in 1971, was the last person raised with the language. Documentation of the Kalapuyan family is mostly of the Central dialect (Santiam), and Yoncalla-specific material is thin. The Wikipedia sources I found say that Kalapuyan revitalization is a goal rather than an established programme. This brief report does not claim otherwise.

\section{Primary Educational Resources}
No school programme, primer, or textbook for Southern Kalapuya was found. The Confederated Tribes of Grand Ronde run a Chinuk Wawa programme, which is a separate language, but it retains Kalapuyan words and place names. I did not find that programme's own pages.

\section{Linguistic and Reference Materials}
\begin{itemize}
\item Wikipedia article on the Kalapuyan languages. \url{https://en.wikipedia.org/wiki/Kalapuyan_languages}
\item Wikipedia article on the Yoncalla language (ISO 639 sxk page). \url{https://wikipedia.com/wiki/ISO_639:sxk}
\item Kalapuyan Dictionary, a University of British Columbia knowledge-base entry. It describes the four-volume dictionary compiled by Paul Stephen McCartney Sr. (Komemma Cultural Protection Association, Yoncalla, 2021), based largely on archival Santiam material from John Hudson. The same source quotes linguists Zenk and Schrock as critical of its linguistic foundations, so it should be used with caution. \url{https://knowledgebase.arts.ubc.ca/?p=8483}
\end{itemize}

\section{Religious and Community Resources}
\begin{itemize}
\item Kalapuya homeland and revitalization pages of the University of Oregon Language Revitalization Lab. They identify descendants as citizens of Grand Ronde, Siletz and other tribes. \url{https://blogs.uoregon.edu/revitalizationlab}
\item Lincoln Chronicle article on Siletz tribal members in southern Oregon who are learning Kalapuyan. \url{https://lincolnchronicle.org/siletz-tribal-members-in-southern-oregon-are-leading-the-effort-to-learn-and-revive-the-native-kalapuyan-language/}
\end{itemize}

\section{Digital and Audio Resources}
\begin{itemize}
\item Portland State University news item on linguist Jedd Schrock, who works on a digital archive of Kalapuya materials for Grand Ronde and on translating Eustace Howard's texts. \url{https://pdx.edu/applied-linguistics/news/applied-linguistics-adjunct-faculty-jedd-schrock-recently-interviewed-reviving-kalapuya}
\item Corvallis Advocate article on Eustace Howard's Kalapuya narrative. \url{https://corvallisadvocate.com/eustace-howard-a-kalapuya-narrative-resurrecte/}
\item Informal blog series, ``Kalapuyan Language Lessons'' (2018), teaching Santiam phrases. Its author is a descendant, not a linguist, and it is Central, not Southern, Kalapuya. \url{https://ndnhistoryresearch.com/?p=5548}
\end{itemize}

\section{Curriculum Development Context}
A heritage curriculum would rest on archival documentation, mainly Melville Jacobs's Kalapuya texts and Santiam materials, so most usable language content is not Southern Kalapuya. A realistic unit would cover the Kalapuya homelands and peoples, place names, and Kalapuyan loanwords in Chinuk Wawa, with Kalapuya words only where scholars have vetted them. Any language work should be coordinated with the Grand Ronde, Siletz, and Yoncalla-area Kalapuya communities and their linguists. Contact details were not found and are not supplied here.
\end{document}